% Part: sets-functions-relations
% Chapter: functions
% Section: kinds

\documentclass[../../../include/open-logic-section]{subfiles}

\begin{document}

\olfileid[tr]{sfr}{fun}{kin}
\olsection{Fonksiyon Türleri}

\begin{explain}
En sık karşılaştığımız bazı fonksiyon türleri için bir sınıflandırma yapmak
yararlı olacaktır.

Önce, değer kümesinin her elemanının fonksiyonun bir değeri olduğu
fonksiyonları ele alabiliriz. Bu tür fonksiyonlara örten denir ve
\olref{fig:surjective}'teki gibi gösterilebilirler.

\begin{figure}
  \olasset{assets/diagrams/surjective.tikz}
  \caption{Örten bir fonksiyonda değer kümesinin her elemanı fonksiyonun bir
    değeridir.}
  \ollabel{fig:surjective}
\end{figure}
\end{explain}

\begin{defn}[Örten fonksiyon]
Bir $f \colon A \rightarrow B$ fonksiyonu, ancak ve ancak $B$ aynı zamanda
$f$'nin görüntü kümesiyse \emph{örtendir}; başka bir deyişle her $y \in B$
için $f(x)=y$ olacak biçimde en az bir $x \in A$ vardır. Simgelerle:
\[
  (\forall y \in B)(\exists x \in A)f(x) = y.
\]
Böyle bir fonksiyona $A$'dan $B$'ye bir örten fonksiyon deriz.
\end{defn}

\begin{explain}
$f$'nin örten olduğunu göstermek istiyorsanız, $f$'nin değer kümesindeki her
nesnenin bir $x$ girdisi için $f(x)$ değeri olduğunu göstermelisiniz.

Her fonksiyonun bir örten fonksiyon \emph{türettiğine} dikkat edin. Gerçekten
$f \colon A \to B$ fonksiyonu verildiğinde, $f'(x)=f(x)$ eşitliğiyle
$f' \colon A \to \ran{f}$ fonksiyonunu tanımlayalım. $\ran{f}$,
\emph{tanım gereği} $\Setabs{f(x)}{x \in A}$ olduğundan, bu $f'$
fonksiyonunun örten olduğu güvence altındadır.
\end{explain}

\begin{explain}
Her fonksiyon olası her girdiyi biricik bir çıktıya eşler. Bunun yanında,
farklı girdileri hiçbir zaman aynı çıktıya eşlemeyen fonksiyonlar da vardır.
Bu tür fonksiyonlara bire bir denir ve
\olref{fig:injective}'teki gibi gösterilebilirler.
\begin{figure}
  \olasset{assets/diagrams/injective.tikz}
  \caption{Bire bir fonksiyon iki farklı argümanı hiçbir zaman aynı
    değere eşlemez.}
  \ollabel{fig:injective}
\end{figure}
\end{explain}

\begin{defn}[Bire bir fonksiyon]
Bir $f \colon A \rightarrow B$ fonksiyonu, ancak ve ancak her $y \in B$ için
$f(x)=y$ olacak biçimde en fazla bir $x \in A$ varsa \emph{bire birdir}.
Böyle bir fonksiyona $A$'dan $B$'ye bir bire bir fonksiyon deriz.
\end{defn}

\begin{explain}
$f$'nin bire bir olduğunu göstermek istiyorsanız, $f$'nin tanım kümesindeki
herhangi iki $x$ ve $y$ elemanı için $f(x)=f(y)$ ise $x=y$ olduğunu
göstermelisiniz.
\end{explain}

\begin{ex}
$f(x)=1$ eşitliğiyle verilen sabit $f\colon \Nat \to \Nat$ fonksiyonu ne
bire birdir ne de örtendir.

$f(x)=x$ eşitliğiyle verilen özdeşlik $f\colon \Nat \to \Nat$ fonksiyonu hem
bire bir hem de örtendir.

$f(x)=x+1$ eşitliğiyle verilen ardıl $f \colon \Nat \to \Nat$ fonksiyonu
bire birdir, fakat örten değildir.

\[
  f(x) =
  \begin{cases}
    \frac{x}{2} & \text{$x$ çiftse} \\
    \frac{x+1}{2} & \text{$x$ tekse.}
  \end{cases}
\]
biçiminde tanımlanan $f \colon \Nat \to \Nat$ fonksiyonu örtendir, fakat
bire bir değildir.
\end{ex}

\begin{explain}
Çoğu zaman hem bire bir hem de örten olan fonksiyonları ele almak isteriz.
Böyle fonksiyonlara bire bir örten denir. Bunlar
\olref{fig:bijective}'te gösterilen fonksiyona benzer. Bire bir örten
fonksiyonlara bazen \emph{bire bir eşlemeler} de denir; çünkü değer kümesinin
elemanlarıyla tanım kümesinin elemanlarını biricik biçimde eşleştirirler.
\begin{figure}
  \olasset{assets/diagrams/bijective.tikz}
  \caption{Bire bir örten bir fonksiyon, değer kümesinin elemanlarıyla tanım
    kümesinin elemanlarını biricik biçimde eşleştirir.}
  \ollabel{fig:bijective}
\end{figure}
\end{explain}

\begin{defn}[Bire bir örten fonksiyon]
Bir $f \colon A \to B$ fonksiyonu, ancak ve ancak hem örten hem de bire bir ise
\emph{bire bir örtendir}. Böyle bir fonksiyona $A$'dan $B$'ye (veya $A$ ile
$B$ arasında) bir bire bir örten fonksiyon deriz.
\end{defn}

\end{document}
